\chapter{Short-time estimates for Gaussian hierarchies}\label{sec:flow}

In this section, we prove the analytic estimates on Gaussian hierarchies that
the three-replica covariance perturbation of \cref{sec:cubic} requires. Near the
window where the covariance path is perturbed, the one-site recursion of that
comparison is a composition of short semilinear parabolic flows on $\R^3$, one
for each of the fixed number of stages of the perturbed path. We must differentiate its value
four times in the width parameter $h$ of the window and twice in the
parameter $k=\delta/h$, uniformly in the starting field. Several covariance
matrices of the path are singular, and the stage of the terminal correction has
a duration proportional to $hk^2$, which vanishes at $k=0$. We state and prove
the estimates for a general finite composition of such flows on $\R^\ell$.

The section has four main results. \Cref{lem:flow-step-limit} identifies the
limit of finite Gaussian hierarchies (\cref{def:inp-hierarchy}) along
refinements with the solution of a flow equation, with an error controlled by
the $L^1$ distance between the masses. \Cref{prop:flow-uniform} bounds all mixed
derivatives of the flow in the parameters $(h,k)$ and in the field, uniformly in
the field, including the endpoint $h=0$ and stages of zero duration, and
\cref{lem:flow-penalty} gives the corresponding bounds for the penalty obtained
by weighting the squared covariance by the mass. \Cref{b:lem:os-prefix} and \cref{cor:flow-prefix}
control the response of a Gaussian hierarchy to a small bounded change of its
terminal function, with an error that does not depend on the number of levels.
In \cref{sec:cubic}, we use the first three results to obtain \cref{b:lem:os-flow}, and
\cref{cor:flow-prefix} is applied to the part of the hierarchy that precedes the
window, in the proof of \cref{b:prop:os-expansion}. From earlier sections, we use
the integration step and the hierarchies of \cref{sec:inputs}
(\cref{def:inp-step,lem:inp-steps,def:inp-hierarchy}). The regularity and
stability of the scalar Parisi equation (\cref{inp:ob-scalar}) enter only the
proof of \cref{cor:flow-prefix}. The remaining tools are the Gaussian
interpolation formula, It\^o's formula, and Gronwall's inequality
\cite[Section~A.13]{talagrand-2011}.

For the scalar Parisi equation, the approximation by step profiles and the
bounds on spatial derivatives that are uniform in the profile are due to
Auffinger and Chen \cite[Propositions~1--2]{auffinger-chen-2015}. The existence,
uniqueness, and spatial smoothness of weak solutions were proved by Jagannath and
Tobasco \cite[Theorem~1]{jagannath-tobasco-2016}. In our
normalization, these facts and the Lipschitz dependence on the profile are
stated in \cref{inp:ob-scalar}, which is \cite[Lemma~3.1]{oai-barriers}. Bovier
and Klimovsky derived the matrix form of the Parisi equation for vector spins
\cite[Section~6, Proposition~6.1 and (6.5)]{bovier-klimovsky-2009}. The
arguments below use the Feynman--Kac representation for the linear equations
obtained by differentiating the flow. We give the estimates for a vector
equation whose diffusion matrices may be singular and whose coefficients depend
on parameters, and we give in detail the bookkeeping of mixed derivatives in the
parameters, which the later sections need.

\paragraph{Notation.}
Throughout this section, $\ell\ge1$ is an integer ($\ell=3$ in \cref{sec:cubic},
one coordinate for each replica), $b>0$ ($b=\beta^2$ in all applications), and
$x\in\R^\ell$ denotes the field variable, with coordinates indexed by
$i,i'\in\{1,\dots,\ell\}$. We write $\langle\cdot,\cdot\rangle$ and
$\abs{\cdot}$ for the Euclidean inner product and norm, $\abs{\Sigma}$ for the
operator norm of a matrix, $U\mathbin{:}V=\sum_{i,i'}U_{ii'}V_{ii'}$ and
$\norm{U}_F^2=U\mathbin{:}U$, and $\Sigma\succeq0$ when $\Sigma$ is symmetric
and positive semidefinite; $\Sigma^{1/2}$ is its positive semidefinite square
root, and $\mathbf g$ is a standard Gaussian vector in $\R^\ell$. For a multi-index
$\iota\in\N^\ell$, $D^\iota$ is the corresponding derivative in $x$ and
$\abs\iota$ its order; $\nabla$ and $D^2$ are the gradient and the Hessian in
$x$. For $\varphi\colon\R^\ell\to\R$ and an integer $r\ge1$, we set
$\norm{D^r\varphi}_\infty=\max_{\abs\iota=r}\sup_x\abs{D^\iota\varphi(x)}$, and
for $2\le r\le\infty$ we let $\mathcal C^r$ be the set of
$\varphi\in C^r(\R^\ell)$ with $\norm{D^{r'}\varphi}_\infty<\infty$ for all finite
$1\le r'\le r$. A function in $\mathcal C^r$ grows at most linearly and need
not itself be bounded. Parameter derivatives are taken in two real variables
$(h,k)$: for $\gamma=(\gamma_h,\gamma_k)\in\N^2$, we write
$\partial^\gamma=\partial_h^{\gamma_h}\partial_k^{\gamma_k}$,
$\abs\gamma=\gamma_h+\gamma_k$, $\gamma'\le\gamma$ for the componentwise order,
$\binom{\gamma}{\gamma'}=\binom{\gamma_h}{\gamma_h'}\binom{\gamma_k}{\gamma_k'}$,
and $\gamma!=\gamma_h!\gamma_k!$. We say that the mixed derivative
$\partial^\gamma D^\iota$ has \emph{weight} $\abs\iota+2\abs\gamma$. One parameter
derivative counts as two spatial derivatives, because a parameter derivative of
a covariance coefficient multiplies a second derivative in $x$.

\section{Flows and the main estimates}\label{sec:flow-statements}

\begin{definition}[Stage data and flows]\label{def:flow}
Let $N\ge1$. A \emph{stage datum} consists of, for each $j\in\{1,\dots,N\}$, a
matrix $B_j\succeq0$, a duration $\lambda_j\ge0$, and a measurable mass function
$m_j\colon[0,1]\to[0,1]$. Stage $j$ occupies the interval $[j-1,j]$ of the
variable $\tau\in[0,N]$. For a terminal function $f\in\mathcal C^2$, the
\emph{flow equation} is
\begin{equation}\label{eq:flow-equation}
 -\partial_\tau F=\frac{b\lambda_j}{2}\Bigl\{B_j\mathbin{:}D^2F
 +m_j(\tau-j+1)\bigl\langle\nabla F,B_j\nabla F\bigr\rangle\Bigr\}
 \quad\text{on }(j-1,j)\times\R^\ell,
\end{equation}
for $j=1,\dots,N$, with $F(N,\cdot)=f$ and $F$ continuous on
$[0,N]\times\R^\ell$. We call $F(0,\cdot)$ the \emph{output} of the flow.
\end{definition}

We solve the equation backward, from $\tau=N$ to $\tau=0$, as the Parisi
equation \eqref{b:eq:Parisi-PDE} is solved backward from time one. Stage~$1$ is
the outermost stage, nearest to the root of the hierarchy, and stage~$N$ is the
innermost. On stage $j$ the Gaussian field has covariance $b\lambda_jB_j$ per unit
of $\tau$; in the notation of \cref{def:inp-hierarchy}, the overlap covariance
path increases by $\lambda_jB_j$ over stage $j$. For example, the Parisi equation
for a distribution function $\alpha$ on an interval $[t_1,t_2]$ is the case
$\ell=1$, $N=1$, $B_1=1$, $\lambda_1=t_2-t_1$, and
$m_1(\tau)=\alpha(t_1+\tau(t_2-t_1))$, with $b=\beta^2$.

\paragraph{Step flows.}
We first define solutions for masses that are step functions. A \emph{step
flow} is the flow of a stage datum in which, for each $j$, there exists a partition
$0=\tau_{j,0}<\tau_{j,1}<\dots<\tau_{j,M_j}=1$ such that $m_j=\mu_{j,i}$ on
$[\tau_{j,i-1},\tau_{j,i})$. With the integration step of
\cref{def:inp-step}, we set $F(N,\cdot)=f$ and
\begin{equation}\label{eq:flow-step-solution}
 F(\tau,\cdot)=T_{\mu_{j,i},\,(j-1+\tau_{j,i}-\tau)b\lambda_jB_j}
 F(j-1+\tau_{j,i},\cdot)
 \quad\text{for } j-1+\tau_{j,i-1}\le\tau\le j-1+\tau_{j,i},
\end{equation}
proceeding from the piece $(N,M_N)$ to the piece $(1,1)$. The formulas for two
consecutive pieces agree at their common endpoint, because
$T_{m,0}\varphi=\varphi$. The output $F(0,x)$ is the value at $x$ of a finite Gaussian hierarchy:
its levels, listed from the root, are the pieces $(j,i)$ in lexicographic order,
and the level $(j,i)$ has covariance $b\lambda_jB_j(\tau_{j,i}-\tau_{j,i-1})$ and
is integrated at mass $\mu_{j,i}$. By \cref{lem:flow-cole-hopf} below, $F$ is
continuous and solves \eqref{eq:flow-equation} on each piece, where the mass is
constant.

If the masses $m_j$ are continuous and $\mathfrak p=(\tau_{j,i})$ is a family of
partitions, the \emph{sampled masses} are $m_j^{\mathfrak p}=m_j(\tau_{j,i-1})$
on $[\tau_{j,i-1},\tau_{j,i})$, and the \emph{mesh} $\abs{\mathfrak p}$ is the
largest length $\tau_{j,i}-\tau_{j,i-1}$. If each $m_j$ is nondecreasing and
$m_j(1)\le m_{j+1}(0)$ for $j<N$, then the sampled masses are nondecreasing from
the root, as in an admissible path of \cref{def:inp-hierarchy}.

\begin{lemma}[Step approximation]\label{lem:flow-step-limit}
Fix a stage datum with matrices $B_j$ and durations $\lambda_j$, and a terminal
function $f\in\mathcal C^2$. Let $L_f=\sup_x\abs{\nabla f(x)}$.
\begin{enumerate}[label=\textup{(\alph*)}]
\item Let $F$ and $\widetilde F$ be the step flows with step masses $m_j$ and
 $\widetilde m_j$. Then $\sup_{\tau,x}\abs{\nabla F(\tau,x)}\le L_f$ and
 \begin{equation}\label{eq:flow-L1}
  \sup_{\tau,x}\abs{F(\tau,x)-\widetilde F(\tau,x)}
  \le\frac{bL_f^2}{2}\sum_{j=1}^N\lambda_j\abs{B_j}
  \int_0^1\abs{m_j(\tau)-\widetilde m_j(\tau)}\dd\tau.
 \end{equation}
 If $F'$ is the step flow with the masses $m_j$ and a terminal function
 $f'\in\mathcal C^2$ such that $f-f'$ is bounded, then
 $\sup_{\tau,x}\abs{F-F'}\le\sup_x\abs{f-f'}$.
\item Let $m_j$ be measurable. For every sequence of step masses $m_j^{(l)}$ with
 $\int_0^1\abs{m_j^{(l)}-m_j}\dd\tau\to0$ for each $j$, the step flows converge
 uniformly on $[0,N]\times\R^\ell$ to a limit $F$ that does not depend on the
 sequence. We call $F$ the \emph{flow} of the stage datum with terminal
 function $f$. The three estimates of \textup{(a)} hold for these flows.
\item Let $1\le N'<N$. Let $F_{\rm in}$ be the flow of stages $N'+1,\dots,N$
 with terminal function $f$, and suppose that $f_{\rm in}=F_{\rm in}(N',\cdot)$
 belongs to $\mathcal C^2$. Then $F$ coincides on $[0,N']\times\R^\ell$ with the
 flow of stages $1,\dots,N'$ with terminal function $f_{\rm in}$.
\item Suppose that $f\in\mathcal C^\infty$ and
 $2\ell b\norm{D^2f}_\infty\sum_j\lambda_j\abs{B_j}\le1$. Then the flow $F$ of
 \textup{(b)} has continuous and bounded spatial derivatives of every order
 $r\ge1$, with $\norm{D^2F(\tau,\cdot)}_\infty\le2\norm{D^2f}_\infty$, and, for
 the step flows $F^{(l)}$ of \textup{(b)} and every $\iota\ne0$,
 \begin{equation}\label{eq:flow-L1-derivatives}
  \sup_{\tau,x}\abs{D^\iota F^{(l)}(\tau,x)-D^\iota F(\tau,x)}
  \le C\sum_{j=1}^N\lambda_j\int_0^1\abs{m^{(l)}_j-m_j}\dd\tau,
 \end{equation}
 where $C$ depends only on $\ell$, $b$, $\abs\iota$, $\max_j\abs{B_j}$,
 $\sum_j\lambda_j\abs{B_j}$, and $\norm{D^{r}f}_\infty$ for $1\le r\le\abs\iota+1$.
 If moreover each $m_j$ is continuous on $[0,1]$, then $F$ is continuously
 differentiable in $\tau$ on each closed stage $[j-1,j]$, with one-sided
 derivatives at its endpoints, and solves \eqref{eq:flow-equation}. It is the
 unique such solution with bounded first and second spatial derivatives that is
 of class $C^{1,2}$ on each closed stage.
\end{enumerate}
\end{lemma}

In words, the hierarchy values converge along refinements of the masses, at a
rate given by the $L^1$ distance between the masses, and the limit is the
solution of the flow equation; for short flows the same holds for every spatial
derivative. Part (c) says that the limiting flows compose stage by stage. Parts
(a)--(c) hold for stages of any duration. The constant in \eqref{eq:flow-L1}
involves only the gradient of $f$. To identify the gradient of the limit in (b),
we also use the bound $K_f+L_f^2$ of \cref{lem:flow-spatial}(a) on the Hessians
of the step flows, where $K_f=\sup_x\abs{D^2f(x)}$. This bound does not require
short stages either. Part (d) is stated only for short flows, which are the only
ones we need.

\paragraph{Parameter-dependent data.}
We now let the data depend on two parameters. Fix $N\ge1$, $h_1>0$, and
$k_0>0$, and let $\mathcal U=[0,h_1]\times[-k_0,k_0]$. A function on $\mathcal U$
is called \emph{smooth} if all its partial derivatives exist and are continuous
on $\mathcal U$; at boundary points, the derivatives are one-sided. In
particular, derivatives in $h$ at $h=0$ are right derivatives. We assume the
following.
\begin{enumerate}[label=\textup{(H\arabic*)}]
\item\label{H1} For each $j$, the map $(h,k)\mapsto B_j(h,k)$ is smooth on
 $\mathcal U$, with values in the positive semidefinite matrices.
\item\label{H2} For each $j$, $\lambda_j(h,k)=h\kappa_j(h,k)$, where
 $\kappa_j\colon\mathcal U\to[0,\infty)$ is smooth. The function $\kappa_j$
 may vanish.
\item\label{H3} For each $j$, $m_j\colon[0,1]\times\mathcal U\to[0,1]$ is
 continuous, and all its derivatives in $(h,k)$ exist and are continuous on
 $[0,1]\times\mathcal U$.
\item\label{H4} The terminal function $f\colon\mathcal U\times\R^\ell\to\R$ is
 smooth in the same sense, and $\partial^\gamma D^\iota f$ is bounded on $\mathcal U\times\R^\ell$
 for every $(\gamma,\iota)\ne(0,0)$.
\item\label{H5} The number $h_0\in(0,h_1]$ satisfies
 $2\ell b\,h_0\Lambda_*\norm{D^2f}_\infty\le1$, where
 $\Lambda_*=\max_{\mathcal U}\sum_j\kappa_j\abs{B_j}$ and
 $\norm{D^2f}_\infty=\sup_{\mathcal U\times\R^\ell}\max_{i,i'}\abs{\partial_{ii'}f}$.
\end{enumerate}
Let $\mathcal U_0=[0,h_0]\times[-k_0,k_0]$. For $(h,k)\in\mathcal U_0$, let
$F(\tau,h,k,x)$ be the flow, in the sense of \cref{lem:flow-step-limit}, of the
stage datum $(B_j(h,k),\lambda_j(h,k),m_j(\cdot,h,k))_{j\le N}$ with terminal
function $f(h,k,\cdot)$, and let $\mathcal B(h,k,x)=F(0,h,k,x)$ be its output.

\begin{proposition}[Uniform bounds for short flows]\label{prop:flow-uniform}
Assume \ref{H1}--\ref{H5}.
\begin{enumerate}[label=\textup{(\roman*)}]
\item For each $(h,k)\in\mathcal U_0$, the function $F(\cdot,h,k,\cdot)$ is
 continuously differentiable in $\tau$ on each closed stage, infinitely
 differentiable in $x$, and it is the unique solution of \eqref{eq:flow-equation}
 with bounded first and second spatial derivatives that is of class $C^{1,2}$ on
 each closed stage.
\item All partial derivatives $\partial^\gamma D^\iota F$ exist and are
 continuous on $[0,N]\times\mathcal U_0\times\R^\ell$. For every
 $(\gamma,\iota)\ne(0,0)$, there exists $C$ such that
 \[
  \sup_{[0,N]\times\mathcal U_0\times\R^\ell}\abs{\partial^\gamma D^\iota F}\le C,
 \]
 where $C$ depends only on $\ell$, $N$, $b$, $h_1$, $\gamma$, and $\abs\iota$,
 on the suprema over $\mathcal U$ of $\abs{\partial^{\gamma'}(\kappa_jB_j)}$ and
 over $[0,1]\times\mathcal U$ of $\abs{\partial^{\gamma'}(\kappa_jm_jB_j)}$, for
 $\gamma'\le\gamma$, and on the suprema over $\mathcal U\times\R^\ell$ of
 $\abs{\partial^{\gamma'}D^{\iota'}f}$ for $\gamma'\le\gamma$ and
 $0<\abs{\iota'}+2\abs{\gamma'}\le\abs\iota+2\abs\gamma$.
 Moreover, $\abs{F(\tau,h,k,x)}\le C_{\rm grow}(1+\abs x)$, where $C_{\rm grow}$
 depends only on $\ell$, $N$, $b$, $h_1$, the supremum over $\mathcal U$ of
 $\abs{\kappa_jB_j}$, the suprema over $\mathcal U\times\R^\ell$ of
 $\abs{\nabla f}$ and $\abs{D^2f}$, and $\sup_{\mathcal U}\abs{f(h,k,0)}$.
\item For $\gamma\ne0$, the derivative $\partial^\gamma F$ is continuously
 differentiable in $\tau$ on each closed stage, and it solves the equation
 obtained by applying $\partial^\gamma$ to \eqref{eq:flow-equation}, with
 terminal value $\partial^\gamma f$.
\item Let $\mathfrak p_l$ be partitions with $\abs{\mathfrak p_l}\to0$, and let
 $F^{\mathfrak p_l}$ be the step flows with the sampled masses $m_j^{\mathfrak p_l}(\cdot,h,k)$.
 Then $F^{\mathfrak p_l}\to F$ uniformly on $[0,N]\times\mathcal U_0\times\R^\ell$, and
 $\partial^\gamma D^\iota F^{\mathfrak p_l}\to\partial^\gamma D^\iota F$ uniformly on compact
 subsets of $[0,N]\times\mathcal U_0\times\R^\ell$, for every $(\gamma,\iota)$.
\end{enumerate}
\end{proposition}

To bound a derivative $\partial^\gamma D^\iota\mathcal B$, we use only
derivatives of the terminal function of weight at most $\abs\iota+2\abs\gamma$
and $(h,k)$-derivatives of order at most $\gamma$ of the coefficients. In
\cref{sec:cubic}, we need the derivatives $\partial_h^{\gamma_h}\partial_k^{\gamma_k}\mathcal B$ with
$\gamma_h\le4$ and $\gamma_k\le2$ together with a few spatial derivatives; with at
most eight spatial derivatives they have weight at most $20$, and their bounds
use $(h,k)$-derivatives of the masses and of the terminal function of order at
most six. The value $F$ itself grows at most linearly in $x$ and is
not estimated, except through its derivatives and through bounded differences
such as $F(\tau,h,k,x)-F(\tau,h',k',x)$.

\begin{remark}[The coefficients at $h=0$]\label{rem:flow-first-order}
Since $\lambda_j(0,k)=0$, we have $F(\tau,0,k,\cdot)=f(0,k,\cdot)$ for all
$\tau$, and hence $\partial_k^{\gamma_k}\mathcal B(0,k,\cdot)=\partial_k^{\gamma_k}f(0,k,\cdot)$ for
every $\gamma_k\ge0$.
Writing $\bar m_j(k)=\int_0^1m_j(\tau,0,k)\dd\tau$ and evaluating $f$, $B_j$ and
$\kappa_j$ at $(0,k)$, we have
\begin{equation}\label{eq:flow-first-order}
 \partial_h\mathcal B(0,k,x)=\partial_hf(0,k,x)+\frac b2\sum_{j=1}^N\kappa_j
 \Bigl\{B_j\mathbin{:}D^2f(0,k,x)
 +\bar m_j\bigl\langle\nabla f(0,k,x),B_j\nabla f(0,k,x)\bigr\rangle\Bigr\}.
\end{equation}
Indeed, at $h=0$ the equation of \cref{prop:flow-uniform}(iii) for
$\gamma=(1,0)$ has no terms involving $\partial_hF$, since the coefficients of
\eqref{eq:flow-equation} vanish, and its source is the derivative in $h$ of the
right side of \eqref{eq:flow-equation} at $F=f(0,k,\cdot)$. Integrating it over
$\tau\in[0,N]$, we obtain \eqref{eq:flow-first-order}.
\end{remark}

\begin{remark}[Singular matrices, zero durations, and regularization]
\label{rem:flow-regularization}
The hypotheses \ref{H1}--\ref{H5} and the constants in \cref{prop:flow-uniform}
involve no lower bound on $B_j$ or on $\kappa_j$. In the proof, every Gaussian
integral is parametrized by its covariance matrix
$t\,bh\kappa_j(h,k)B_j(h,k)$, which is a smooth function of $(h,k)$ with values
in the positive semidefinite matrices, including where it vanishes, and
\cref{lem:flow-gauss} differentiates Gaussian integrals as functions of the
covariance. No square root of a covariance or of a duration is differentiated.
For instance, a stage with $\kappa_j(h,k)=ck^2$ has zero duration at $k=0$, and
its contribution is differentiable in $k$ there through the polynomial $hk^2$.
One may also replace each $B_j$ by $B_j+\eps\mathrm{Id}$ with $0<\eps\le1$. Let
$K_*=\max_{\mathcal U}\sum_j\kappa_j$, and decrease $h_0$, if necessary, so that
$2\ell b\,h_0(\Lambda_*+K_*)\norm{D^2f}_\infty\le1$. The regularized data then
satisfy \ref{H1}--\ref{H5}, with upper bounds uniform in $\eps$. For parameters
with $h\kappa_j(h,k)>0$, the diffusion matrix of stage $j$ is positive definite.
Stages with $h\kappa_j(h,k)=0$, in particular all stages at $h=0$, still have zero
diffusion, and the regularized equations are not uniformly parabolic on
$\mathcal U_0$. The
regularized flows satisfy the bounds of \cref{prop:flow-uniform} uniformly in
$\eps$, and they converge to $F$ as $\eps\to0$, together with all derivatives
$\partial^\gamma D^\iota$, uniformly on compact sets. Indeed, the difference
between the regularized flow and $F$ satisfies the linear equation in the proof
of \cref{lem:flow-step-limit}(a), with $B_j+\eps\mathrm{Id}$ in place of $B_j$ and with
equal masses. Its source is
$(b\lambda_j\eps/2)\{\mathrm{Id}\mathbin{:}D^2F+m_j\abs{\nabla F}^2\}$, which is bounded
by $C\eps$. Then \cref{lem:flow-linear} shows that the values
converge, and the compactness argument in the proof of
\cref{prop:flow-uniform}(iv) identifies the limits of all derivatives. We do not
use the regularization.
\end{remark}

\begin{remark}[Further parameters and restriction to a line]
\label{rem:flow-uniformity}
If the data depend on a further parameter in which no derivatives are needed
(in \cref{sec:cubic}, the left endpoint of the window), and the suprema listed in
\cref{prop:flow-uniform}(ii) and the number $\Lambda_*\norm{D^2f}_\infty$ are
bounded uniformly in it, then the conclusions hold uniformly in that parameter;
for the linear growth bound, $\sup_{\mathcal U}\abs{f(h,k,0)}$ must also be
bounded uniformly in that parameter.
The statements and the proofs of this section hold without change if $k$ is a
vector of parameters ranging over a product of compact intervals, with
multi-indices in place of $\gamma_k$; a further parameter in which derivatives
are needed can be included in this way. For a vector $e\in\R^\ell$ and $z\in\R$,
the restriction $z\mapsto\mathcal B(h,k,ze)$
satisfies
\[
 \partial^\gamma\partial_z^r\bigl[\mathcal B(h,k,ze)\bigr]
 =\sum_{\abs\iota=r}\frac{r!}{\iota!}\,e^\iota\,\partial^\gamma D^\iota\mathcal B(h,k,ze),
\]
where $\iota!=\prod_i\iota_i!$ and $e^\iota=\prod_ie_i^{\iota_i}$. Its derivatives
obey the bounds of \cref{prop:flow-uniform}(ii), multiplied by
$(\sum_i\abs{e_i})^r$.
\end{remark}

The last estimate of this section concerns the penalty. For a smooth function
$Q_\star\colon\mathcal U\to\R^{\ell\times\ell}$ with symmetric values, let
\[
 Q(\tau)=Q_\star+\sum_{j'<j}\lambda_{j'}B_{j'}+(\tau-j+1)\lambda_jB_j
 \qquad(j-1\le\tau\le j),
\]
where all functions are evaluated at $(h,k)$. This is the overlap covariance
path through the stages, started at $Q_\star$. We set
\begin{equation}\label{eq:flow-penalty}
 \Pi(h,k)=\sum_{j=1}^N\int_{j-1}^jm_j(\tau-j+1)\,\partial_\tau\norm{Q(\tau)}_F^2\dd\tau
 =2\sum_{j=1}^N\lambda_j\int_0^1m_j(\tau)\,Q(j-1+\tau)\mathbin{:}B_j\dd\tau.
\end{equation}
This is the contribution of the stages to the penalty $\int m\dd\norm{Q}_F^2$
of \eqref{eq:inp-penalty}. For partitions $\mathfrak p$, the corresponding
contribution of the step hierarchy with the sampled masses is
\[
 \Pi^{\mathfrak p}(h,k)=\sum_{j=1}^N\sum_{i=1}^{M_j}m_j(\tau_{j,i-1})
 \Bigl(\norm{Q(j-1+\tau_{j,i})}_F^2-\norm{Q(j-1+\tau_{j,i-1})}_F^2\Bigr).
\]

\begin{lemma}[The penalty]\label{lem:flow-penalty}
Assume \ref{H1}--\ref{H3}. Then $\Pi$ is smooth on $\mathcal U$, and for every
$\gamma$ we have $\sup_{\mathcal U}\abs{\partial^\gamma\Pi}\le C$, where $C$
depends only on $\ell$, $N$, $h_1$, $\gamma$, and the suprema over $\gamma'\le\gamma$
of $\abs{\partial^{\gamma'}Q_\star}$, $\abs{\partial^{\gamma'}(\kappa_jB_j)}$ and
$\abs{\partial^{\gamma'}(\kappa_jm_jB_j)}$. Moreover
$\sup_{\mathcal U}\abs{\partial^\gamma(\Pi^{\mathfrak p}-\Pi)}\to0$ as
$\abs{\mathfrak p}\to0$, for every $\gamma$.
\end{lemma}

\begin{proof}
Since $m_j^{\mathfrak p}$ is constant on each piece, we have
$\Pi^{\mathfrak p}=2\sum_j\lambda_j\int_0^1m_j^{\mathfrak p}(\tau)\,Q(j-1+\tau)\mathbin{:}B_j\dd\tau$,
which is \eqref{eq:flow-penalty} with $m_j$ replaced by $m_j^{\mathfrak p}$. The
integrand on stage $j$ is $Q(j-1+\tau)\mathbin{:}(h\kappa_jm_jB_j)$, where
$Q(j-1+\tau)$ is a polynomial in $\tau$, $Q_\star$, and the matrices
$h\kappa_{j'}B_{j'}$, and the $(h,k)$-derivatives of $m_j^{\mathfrak p}$ are those
of $m_j$ evaluated at the left endpoints of the pieces. Differentiating under the integral sign, we
find that $\Pi$ and $\Pi^{\mathfrak p}$ are smooth with the stated bounds. By
\ref{H3}, every $(h,k)$-derivative of $m_j$ is uniformly continuous on the
compact set $[0,1]\times\mathcal U$, and we have
$\partial^{\gamma'}m_j^{\mathfrak p}\to\partial^{\gamma'}m_j$ uniformly as
$\abs{\mathfrak p}\to0$. This implies the last statement.
\end{proof}

\section{Three tools}\label{sec:flow-tools}

With the first tool, we differentiate a Gaussian integral with respect to its
covariance, including at singular covariances.

\begin{lemma}[Gaussian integrals as functions of the covariance]\label{lem:flow-gauss}
Let $\Omega$ be a convex subset of a Euclidean space. Let $\mathcal F\colon\Omega\times\R^\ell\to\R$ be
continuous and twice continuously differentiable in its second variable $y$,
with $D^2_y\mathcal F$ continuous on $\Omega\times\R^\ell$, and suppose that there exists
$c<\infty$ such that
\begin{equation}\label{eq:flow-growth}
 \abs{\mathcal F(\omega,y)}+\abs{\nabla_y\mathcal F(\omega,y)}+\abs{D^2_y\mathcal F(\omega,y)}
 \le c\,e^{c\abs y}\qquad(\omega\in\Omega,\ y\in\R^\ell).
\end{equation}
\begin{enumerate}[label=\textup{(\alph*)}]
\item The map $(\omega,\Sigma)\mapsto\E\mathcal F(\omega,\Sigma^{1/2}\mathbf g)$ is continuous
 on $\Omega\times\{\Sigma\succeq0\}$. This holds for every continuous $\mathcal F$ with
 $\abs{\mathcal F}\le ce^{c\abs y}$.
\item For $\omega\in\Omega$ and $\Sigma,\Sigma'\succeq0$, with
 $\Sigma_t=\Sigma+t(\Sigma'-\Sigma)$,
 \[
  \E\mathcal F(\omega,\Sigma'^{1/2}\mathbf g)-\E\mathcal F(\omega,\Sigma^{1/2}\mathbf g)
  =\frac12\int_0^1(\Sigma'-\Sigma)\mathbin{:}
  \E D^2_y\mathcal F(\omega,\Sigma_t^{1/2}\mathbf g)\dd t.
 \]
\item Let $\omega\mapsto\Sigma(\omega)\succeq0$ be continuously differentiable on
 $\Omega$, and suppose that $\partial_{\omega_a}\mathcal F$ exists, is continuous, and
 satisfies $\abs{\partial_{\omega_a}\mathcal F}\le ce^{c\abs y}$. Then the function
 $\omega\mapsto\E\mathcal F(\omega,\Sigma(\omega)^{1/2}\mathbf g)$ has a continuous partial
 derivative in $\omega_a$, one-sided at boundary points of $\Omega$, and
 \begin{equation}\label{eq:flow-gauss-derivative}
  \partial_{\omega_a}\E\mathcal F(\omega,\Sigma(\omega)^{1/2}\mathbf g)
  =\E\partial_{\omega_a}\mathcal F(\omega,\Sigma(\omega)^{1/2}\mathbf g)
  +\frac12\,\partial_{\omega_a}\Sigma(\omega)\mathbin{:}
  \E D^2_y\mathcal F(\omega,\Sigma(\omega)^{1/2}\mathbf g).
 \end{equation}
\end{enumerate}
\end{lemma}

\begin{proof}
For (a), recall that the positive square root is continuous on the positive
semidefinite matrices. Indeed, if the spectra of $\Sigma$ and $\Sigma'$ lie in
$[0,c']$, then $\abs{\Sigma^{1/2}-\Sigma'^{1/2}}\le2\sup_{[0,c']}\abs{\sqrt s-P(s)}
+\abs{P(\Sigma)-P(\Sigma')}$ for every polynomial $P$, and $\sqrt s$ is a uniform
limit of polynomials on $[0,c']$. If $(\omega_n,\Sigma_n)\to(\omega,\Sigma)$, then
$\mathcal F(\omega_n,\Sigma_n^{1/2}\mathbf g)\to\mathcal F(\omega,\Sigma^{1/2}\mathbf g)$ almost surely, and
these random variables are dominated by $c\exp(cc'^{1/2}\abs{\mathbf g})$, where $c'$ bounds
the spectra of all $\Sigma_n$. By dominated convergence, the expectations
converge, which proves (a).

For (b), let $\mathbf g'$ be a standard Gaussian vector independent of $\mathbf g$, and let
$\zeta(t)=\sqrt t\,\Sigma'^{1/2}\mathbf g'+\sqrt{1-t}\,\Sigma^{1/2}\mathbf g$. The vector $\zeta(t)$ is
centered Gaussian with covariance $\Sigma_t$, so it has the law of
$\Sigma_t^{1/2}\mathbf g$. By the Gaussian interpolation formula
\cite[Lemma~1.3.1]{talagrand-2011}, for $0<t<1$,
\[
 \frac{\dd}{\dd t}\E\mathcal F(\omega,\zeta(t))
 =\frac12(\Sigma'-\Sigma)\mathbin{:}\E D_y^2\mathcal F(\omega,\zeta(t)).
\]
That lemma is stated for infinitely differentiable functions whose derivatives
grow more slowly than $e^{a\abs y^2}$ for every $a>0$. For $\mathcal F(\omega,\cdot)$
as in \eqref{eq:flow-growth}, we apply it to the convolution of $\mathcal F(\omega,\cdot)$
with a smooth compactly supported mollifier, integrate the identity between two
times in $(0,1)$, and let the mollifier shrink; both sides converge by dominated
convergence. The right side of the displayed identity is continuous on $[0,1]$ by
(a), applied to $D^2_y\mathcal F$. Integrating over $t\in[0,1]$, we obtain (b).

For (c), fix $\omega$ and write $\Sigma=\Sigma(\omega)$ and
$\Sigma_\eta=\Sigma(\omega+\eta e_a)$ for $\omega+\eta e_a\in\Omega$, where
$e_a$ is the $a$th coordinate vector. Then
\begin{multline*}
 \E\mathcal F(\omega+\eta e_a,\Sigma_\eta^{1/2}\mathbf g)-\E\mathcal F(\omega,\Sigma^{1/2}\mathbf g)
 =\int_0^\eta\E\partial_{\omega_a}\mathcal F(\omega+\eta'e_a,\Sigma_\eta^{1/2}\mathbf g)\dd\eta'\\
 +\frac12\int_0^1(\Sigma_\eta-\Sigma)\mathbin{:}
 \E D_y^2\mathcal F\bigl(\omega,(\Sigma+t(\Sigma_\eta-\Sigma))^{1/2}\mathbf g\bigr)\dd t,
\end{multline*}
where we used the fundamental theorem of calculus and Fubini's theorem for the
first term and (b) for the second; the segment from $\Sigma$ to $\Sigma_\eta$
consists of positive semidefinite matrices. We divide by $\eta$ and let
$\eta\to0$. By (a), applied to $\partial_{\omega_a}\mathcal F$ and to the entries of
$D^2_y\mathcal F$, the two integrands converge uniformly to their values at $\eta=0$,
and we obtain \eqref{eq:flow-gauss-derivative}. The right side of
\eqref{eq:flow-gauss-derivative} is continuous by (a).
\end{proof}

The second tool collects the properties of one integration step
$T_{m,\Sigma}$ of \cref{def:inp-step} that we need. Recall that monotonicity, the
translation identity, and the semigroup property are \cref{lem:inp-steps}(a)
and (b). For constant mass, the exponential change of variables that turns
the flow into the heat equation is the Cole--Hopf transformation; its scalar
Parisi form is recorded in \cite[Section~2, (10)--(12)]{auffinger-chen-2015}.
The next lemma gives the corresponding finite-dimensional Gaussian step.

\begin{lemma}[One integration step]\label{lem:flow-cole-hopf}
Let $\Sigma\succeq0$, let $m\in[0,1]$, and let $\varphi\in\mathcal C^2$.
\begin{enumerate}[label=\textup{(\alph*)}]
\item We have $T_{m,\Sigma}\varphi\in\mathcal C^2$ and
 $\sup_x\abs{\nabla T_{m,\Sigma}\varphi(x)}\le\sup_x\abs{\nabla\varphi(x)}$. If
 $\varphi\in\mathcal C^r$, then $T_{m,\Sigma}\varphi\in\mathcal C^r$.
\item The function $\psi(t,x)=T_{m,t\Sigma}\varphi(x)$, $t\ge0$, is continuous and
 continuously differentiable in $t$, with a right derivative at $t=0$; its first
 and second derivatives in $x$ are continuous in $(t,x)$; and
 \begin{equation}\label{eq:flow-one-step-pde}
  \partial_t\psi=\frac12\Bigl\{\Sigma\mathbin{:}D^2\psi
  +m\bigl\langle\nabla\psi,\Sigma\nabla\psi\bigr\rangle\Bigr\}.
 \end{equation}
\item Let $\Omega$ be a compact convex subset of a Euclidean space, and let $\Sigma(\omega)\succeq0$,
 $m(\omega)\in[0,1]$, and $\varphi(\omega,x)$ be smooth functions of
 $\omega\in\Omega$ and of $(\omega,x)\in\Omega\times\R^\ell$, such that
 $\partial^\gamma_\omega D^\iota\varphi$ is bounded for every $(\gamma,\iota)\ne(0,0)$.
 Then $\psi(t,\omega,x)=T_{m(\omega),t\Sigma(\omega)}\varphi(\omega,\cdot)(x)$ is
 infinitely differentiable in $(t,\omega,x)\in[0,1]\times\Omega\times\R^\ell$
 (one-sided at the boundary), the function $\psi-\varphi$ is bounded, and every
 derivative $\partial^\gamma_\omega D^\iota\psi$ with $(\gamma,\iota)\ne(0,0)$ is
 bounded on $[0,1]\times\Omega\times\R^\ell$.
\end{enumerate}
\end{lemma}

\begin{proof}
Let $L_\varphi=\sup_x\abs{\nabla\varphi(x)}$, and, for fixed $x$, let
$\mathcal Y_x=\varphi(x+\Sigma^{1/2}\mathbf g)-\varphi(x)$, so that $\abs{\mathcal Y_x}\le L_\varphi\abs{\Sigma^{1/2}\mathbf g}$.
For $m>0$, we have
$T_{m,\Sigma}\varphi(x)=\varphi(x)+m^{-1}\log\E e^{m\mathcal Y_x}$.

For (a), since the derivatives of $x\mapsto e^{m\varphi(x+y)-m\varphi(x)}$ of
order at most two are bounded by $C(1+\abs y)^2e^{L_\varphi\abs y}$, we may differentiate
under the expectation. We obtain
\begin{equation}\label{eq:flow-tilted-gradient}
 \nabla T_{m,\Sigma}\varphi(x)
 =\frac{\E\bigl[\nabla\varphi(x+\Sigma^{1/2}\mathbf g)\,e^{m\mathcal Y_x}\bigr]}{\E e^{m\mathcal Y_x}}.
\end{equation}
The right side is an average of values of $\nabla\varphi$ with respect to a
probability measure, so its norm is at most $L_\varphi$.
Differentiating further, we find that $D^\iota T_{m,\Sigma}\varphi(x)$, for
$1\le\abs\iota\le r$, is a polynomial, with coefficients depending only on $\iota$
and $m$, in the averages
$\E[P(x+\Sigma^{1/2}\mathbf g)e^{m\mathcal Y_x}]/\E e^{m\mathcal Y_x}$ of products $P$ of the derivatives
$D^{\iota'}\varphi$ with $1\le\abs{\iota'}\le\abs\iota$. These derivatives are
bounded, and so are the averages. The case $m=0$ is the same without the
exponential weights.

For (b) and $m>0$, let $U(t,x)=\E\exp\{m\varphi(x+(t\Sigma)^{1/2}\mathbf g)\}$. The
function $y\mapsto e^{m\varphi(x+y)}$ and its first two derivatives are bounded
by $Ce^{m\abs{\varphi(x)}+L_\varphi\abs y}$, locally uniformly in $x$. By
\cref{lem:flow-gauss}(a) and (b), with the covariances $t\Sigma$ and
$t'\Sigma$, the function $U$ is continuous, continuously differentiable in $t$,
and $\partial_tU=(1/2)\Sigma\mathbin{:}D^2U$. By \cref{lem:flow-gauss}(a),
$\nabla U$ and $D^2U$ are continuous in $(t,x)$. Writing $\psi=m^{-1}\log U$, we
have $\nabla U=mU\nabla\psi$ and $D^2U=mU\{D^2\psi+m\nabla\psi\nabla\psi^{\mathsf T}\}$,
and \eqref{eq:flow-one-step-pde} follows. For $m=0$, we apply
\cref{lem:flow-gauss}(b) to $\varphi$ directly.

For (c), we write $m=m(\omega)$, $\Sigma=\Sigma(\omega)$, and
$\mathcal Y=\varphi(\omega,x+(t\Sigma)^{1/2}\mathbf g)-\varphi(\omega,x)$. For every $m\ge0$,
\begin{equation}\label{eq:flow-ch-integral}
 \psi(t,\omega,x)-\varphi(\omega,x)
 =\int_0^1\frac{\E\bigl[\mathcal Y\,e^{\rho m\mathcal Y}\bigr]}{\E e^{\rho m\mathcal Y}}\dd\rho.
\end{equation}
Indeed, for $m>0$ the integrand is the derivative of
$\rho\mapsto m^{-1}\log\E e^{\rho m\mathcal Y}$, and for $m=0$ both sides equal $\E \mathcal Y$.
We write $\widehat\varphi(\omega,x,y)=\varphi(\omega,x+y)-\varphi(\omega,x)$. The
denominator is the Gaussian integral $\E\mathcal F(\omega,x,\rho;(t\Sigma)^{1/2}\mathbf g)$ of
the function
\[
 \mathcal F(\omega,x,\rho;y)=\exp\bigl\{\rho\,m(\omega)\,\widehat\varphi(\omega,x,y)\bigr\},
\]
and the numerator is the Gaussian integral of $\widehat\varphi\,\mathcal F$. Every
derivative of $\mathcal F$ and of $\widehat\varphi\,\mathcal F$ in $(\omega,x,\rho,y)$ is a
polynomial in
$\varphi(\omega,x+y)-\varphi(\omega,x)$, in the derivatives of $m$, and in the
derivatives $\partial^{\gamma'}_\omega D^{\iota'}\varphi$ with
$(\gamma',\iota')\ne(0,0)$ evaluated at $(\omega,x+y)$ and $(\omega,x)$, times
the exponential factor. Since $\abs{\varphi(\omega,x+y)-\varphi(\omega,x)}\le
L_\varphi\abs y$ with $L_\varphi=\sup\abs{\nabla_x\varphi}$, all these derivatives are bounded by
$C(1+\abs y)^{C}e^{L_\varphi\abs y}$, uniformly in $(t,\omega,x,\rho)$. The covariance
$(t,\omega)\mapsto t\Sigma(\omega)$ is smooth with positive semidefinite values.
Applying \cref{lem:flow-gauss}(c) repeatedly, with the parameter
$(t,\omega,x,\rho)$ in the convex set $[0,1]\times\Omega\times\R^\ell\times[0,1]$,
we find that the numerator and the denominator are infinitely differentiable,
and that each of their derivatives is a Gaussian integral of a function with the
same bound; all these derivatives are bounded. The denominator is
at least $\E e^{-L_\varphi\abs{(t\Sigma)^{1/2}\mathbf g}}$, which is bounded below by a positive
constant uniformly in $t\le1$ and $\omega\in\Omega$, since $\Omega$ is compact.
Then the right side of \eqref{eq:flow-ch-integral} is infinitely
differentiable, with all its derivatives and the function itself bounded. The
derivatives of $\varphi$ of order at least one are bounded by hypothesis, and
this proves (c).
\end{proof}

The third tool is a bound for linear equations of the type
\eqref{eq:flow-equation}. Its proof uses the stochastic representation of the
solution and needs no lower bound on the diffusion matrix. If the
diffusion matrix is degenerate, the representing diffusion moves in fewer
directions.

\begin{lemma}[A linear estimate]\label{lem:flow-linear}
Let $T_*>0$ and suppose that $[0,T_*]$ is the union of finitely many closed
intervals with disjoint interiors. Let $\mathcal A(t)\succeq0$ be constant on each
of them. Let $\phi\colon[0,T_*]\times\R^\ell\to\R^\ell$ and
$\Upsilon\colon[0,T_*]\times\R^\ell\to\R$ be bounded and measurable, with $\phi$
Lipschitz in $x$ uniformly in $t$, and with $t\mapsto\sup_x\abs{\Upsilon(t,x)}$
measurable. Let $\psi\colon[0,T_*]\times\R^\ell\to\R$ be continuous, with
$\abs{\psi(t,x)}+\abs{\nabla\psi(t,x)}\le c(1+\abs x)$, and suppose that on each
of the intervals $\psi$ coincides with a function of class $C^{1,2}$ and satisfies
\[
 -\partial_t\psi=\mathcal A(t)\mathbin{:}D^2\psi+\langle\phi,\nabla\psi\rangle+\Upsilon
\]
in its interior. Then, for every $t\in[0,T_*]$,
\begin{equation}\label{eq:flow-linear}
 \sup_x\abs{\psi(t,x)}\le\sup_x\abs{\psi(T_*,x)}
 +\int_t^{T_*}\sup_x\abs{\Upsilon(t',x)}\dd t'.
\end{equation}
\end{lemma}

\begin{proof}
Fix $t$ and $x$, let $W$ be a standard Brownian motion in $\R^\ell$, and let
$\mathcal X$ solve
\[
 \mathcal X_{t'}=x+\int_t^{t'}\phi(t'',\mathcal X_{t''})\dd t''
 +\int_t^{t'}\bigl(2\mathcal A(t'')\bigr)^{1/2}\dd W_{t''},
 \qquad t\le t'\le T_*.
\]
Since the noise term does not depend on $\mathcal X$ and $\phi$ is bounded and Lipschitz
in space, this equation has a unique continuous adapted solution, obtained
pathwise by Picard iteration and Gronwall's inequality. On each of the finitely
many intervals, It\^o's formula \cite[Theorem~3.3.6]{karatzas-shreve-1991} applies
to the $C^{1,2}$ function that coincides with $\psi$ there. Adding the results and
using the equation for $\psi$, we obtain
\[
 \psi(T_*,\mathcal X_{T_*})-\psi(t,x)
 =-\int_t^{T_*}\Upsilon(t',\mathcal X_{t'})\dd t'
 +\int_t^{T_*}\bigl\langle\nabla\psi(t',\mathcal X_{t'}),
 \bigl(2\mathcal A(t')\bigr)^{1/2}\dd W_{t'}\bigr\rangle.
\]
We have $\E\sup_{t'}\abs{\mathcal X_{t'}}^2<\infty$ because $\phi$ is bounded, and the
stochastic integral has mean zero because $\nabla\psi$ has linear growth. Taking expectations, we
find
\[
 \psi(t,x)=\E\psi(T_*,\mathcal X_{T_*})+\E\int_t^{T_*}\Upsilon(t',\mathcal X_{t'})\dd t',
\]
and \eqref{eq:flow-linear} follows. This is the Feynman--Kac representation
\cite[Theorem~5.7.6]{karatzas-shreve-1991}, which we derived directly because
$\psi$ is of class $C^{1,2}$ only on each of the intervals.
\end{proof}

\section{Spatial derivatives and the step limit}\label{sec:flow-spatial}

We first bound the spatial derivatives of step flows. The constants depend only
on the durations and the matrices, not on the masses or the partitions.

\begin{lemma}[Spatial bounds for step flows]\label{lem:flow-spatial}
Let $F$ be a step flow with terminal function $f$.
\begin{enumerate}[label=\textup{(\alph*)}]
\item If $f\in\mathcal C^2$, then, for every $\tau$,
 \[
  \sup_x\abs{\nabla F(\tau,x)}\le L_f,\qquad
  \sup_x\abs{D^2F(\tau,x)}\le K_f+L_f^2,
 \]
 where $L_f=\sup_x\abs{\nabla f(x)}$, $K_f=\sup_x\abs{D^2f(x)}$, and $\abs{\cdot}$
 is the operator norm. This holds for stages of any duration and any masses in
 $[0,1]$.
\item If $f\in\mathcal C^\infty$ and
 $2\ell b\norm{D^2f}_\infty\sum_j\lambda_j\abs{B_j}\le1$, then
 $\norm{D^2F(\tau,\cdot)}_\infty\le2\norm{D^2f}_\infty$ for every $\tau$.
\item Under the hypotheses of \textup{(b)}, for every $r\ge3$ there exists $C$,
 depending only on $\ell$, $b$, $r$, $\sum_j\lambda_j\abs{B_j}$, and
 $\norm{D^{r'}f}_\infty$ for $2\le r'\le r$, such that
 $\norm{D^rF(\tau,\cdot)}_\infty\le C$ for every $\tau$.
\end{enumerate}
\end{lemma}

\begin{proof}
For (a), we use the weighted martingale-variance decomposition in the proof of
\cite[Lemma~3.1]{oai-barriers}, with a general vector terminal function in place
of the scalar terminal function $\log(2\cosh z)$.
The gradient bound in (a) follows from \cref{lem:flow-cole-hopf}(a), applied to
each step of \eqref{eq:flow-step-solution}. For the Hessian, fix $(\tau,x)$ and a
unit vector $v$. By \eqref{eq:flow-step-solution}, $F(\tau,x)$ is the root value
of a finite hierarchy, whose levels are the pieces after $\tau$; if $\tau$ lies
inside a piece, the first level is the remaining part of that piece. We list the
levels from the root as $1,\dots,M$, with covariances $\Sigma_j$ and masses
$\mu_j\in[0,1]$, and let $\Psi_M=f$ and $\Psi_{j-1}=T_{\mu_j,\Sigma_j}\Psi_j$, so
that $\Psi_0=F(\tau,\cdot)$. Each $\Psi_j$ belongs to $\mathcal C^2$ by
\cref{lem:flow-cole-hopf}(a). Let $z_0=x$, and given $z_0,\dots,z_{j-1}$, let $z_j$
have the law $\mathcal K_{\mu_j,\Sigma_j,\Psi_j}(z_{j-1},\cdot)$ of
\eqref{eq:inp-kernel}. The weights in \eqref{eq:flow-tilted-gradient} are the
densities of this kernel, and therefore
$\partial_v\Psi_{j-1}(z)=\E[\partial_v\Psi_j(z_j)\mid z_{j-1}=z]$.
Differentiating \eqref{eq:flow-tilted-gradient} once more in the direction $v$,
under the expectation as in the proof of \cref{lem:flow-cole-hopf}(a), we find
that
\[
 \partial_v^2\Psi_{j-1}(z)=\E\bigl[\partial_v^2\Psi_j(z_j)\bigm|z_{j-1}=z\bigr]
 +\mu_j\Var\bigl(\partial_v\Psi_j(z_j)\bigm|z_{j-1}=z\bigr).
\]
The process $\mathcal M_j=\partial_v\Psi_j(z_j)$ is a bounded martingale, by the
Markov property, and the conditional variance in the display, evaluated at
$z=z_{j-1}$, has expectation $\E(\mathcal M_j-\mathcal M_{j-1})^2$. Iterating the
display from $j=1$ to $j=M$, we obtain
\[
 \partial_v^2F(\tau,x)=\E\,\partial_v^2f(z_M)
 +\sum_{j=1}^M\mu_j\E\bigl(\mathcal M_j-\mathcal M_{j-1}\bigr)^2.
\]
Since $0\le\mu_j\le1$ and the increments of a martingale are orthogonal, the sum
lies between $0$ and $\sum_j\E(\mathcal M_j-\mathcal M_{j-1})^2=\Var(\mathcal M_M)\le
L_f^2$. We obtain $-K_f\le\partial_v^2F(\tau,x)\le K_f+L_f^2$ for every unit vector
$v$, and since $D^2F(\tau,x)$ is symmetric, this is the Hessian bound. The bound
does not depend on $M$, on the covariances, or on the order of the masses. This
is the argument of the proof of \cref{b:lem:os-prefix} below, with a derivative
in the field in place of a derivative in the terminal perturbation.

For (b) and (c), we note that, by \cref{lem:flow-cole-hopf}(c) with a one-point
parameter set, $F(\tau,\cdot)\in\mathcal C^\infty$ for every $\tau$, all spatial
derivatives of $F$ are continuous, and on each closed piece they are continuously
differentiable in $\tau$ and satisfy the equations obtained by differentiating
\eqref{eq:flow-equation} in $x$. On each piece, each spatial derivative of
$F$ is bounded. For $r\ge1$, let
$\Theta_r(\tau)=\max_{\abs\iota=r}\sup_x\abs{D^\iota F(\tau,x)}$. It is bounded on
$[0,N]$ and lower semicontinuous, hence measurable. On stage $j$ we write
$\lambda(\tau)\abs{B(\tau)}=\lambda_j\abs{B_j}$.

Differentiating \eqref{eq:flow-equation} in $x_i$ and $x_{i'}$, we find that
$\psi=\partial_{ii'}F$ satisfies, on stage $j$,
\[
 -\partial_\tau\psi=\frac{b\lambda_j}{2}B_j\mathbin{:}D^2\psi
 +b\lambda_jm_j\bigl\langle B_j\nabla F,\nabla\psi\bigr\rangle
 +b\lambda_jm_j\bigl\langle\nabla\partial_iF,B_j\nabla\partial_{i'}F\bigr\rangle.
\]
We apply \cref{lem:flow-linear} with $\mathcal A=(b\lambda_j/2)B_j$, with the drift
$\phi=b\lambda_jm_jB_j\nabla F$, which is bounded by (a) and Lipschitz in $x$
since $D^2F$ is bounded on each piece, and with the last term as $\Upsilon$.
Using $0\le m_j\le1$ and $\abs{\nabla\partial_iF}\le\sqrt\ell\,\Theta_2$, we find
that $\abs\Upsilon\le\ell b\lambda_j\abs{B_j}\Theta_2^2$, and \cref{lem:flow-linear}
yields
\[
 \Theta_2(\tau)\le \Xi(\tau),\qquad
 \Xi(\tau)=\norm{D^2f}_\infty+\ell b\int_\tau^N\lambda(\tau')\abs{B(\tau')}\,
 \Theta_2(\tau')^2\dd\tau'.
\]
If $\norm{D^2f}_\infty=0$, then $\Xi(\tau)\le C\int_\tau^N\Xi(\tau')\dd\tau'$, and
Gronwall's inequality implies that $\Xi=0$. Otherwise $\Xi$ is positive and absolutely
continuous, and $-\Xi'=\ell b\lambda\abs B\Theta_2^2\le\ell b\lambda\abs B\,\Xi^2$
almost everywhere. Then $(1/\Xi)'\le\ell b\lambda\abs B$, and integrating from
$\tau$ to $N$, we obtain
\[
 \frac1{\Xi(\tau)}\ge\frac1{\norm{D^2f}_\infty}-\ell b\sum_j\lambda_j\abs{B_j}
 \ge\frac1{2\norm{D^2f}_\infty},
\]
where we used the hypothesis of (b) in the last step. This proves (b).

For (c), let $\abs\iota=r\ge3$. By the Leibniz rule,
\begin{multline*}
 -\partial_\tau D^\iota F=\frac{b\lambda_j}{2}B_j\mathbin{:}D^2D^\iota F
 +b\lambda_jm_j\bigl\langle B_j\nabla F,\nabla D^\iota F\bigr\rangle\\
 +\frac{b\lambda_jm_j}{2}\sum_{0<\iota'<\iota}\binom{\iota}{\iota'}
 \bigl\langle\nabla D^{\iota'}F,B_j\nabla D^{\iota-\iota'}F\bigr\rangle,
\end{multline*}
where $0<\iota'<\iota$ means $\iota'\le\iota$ componentwise with
$\iota'\notin\{0,\iota\}$. The two factors in the last sum have orders
$\abs{\iota'}+1$ and $r-\abs{\iota'}+1$, both between $2$ and $r$. In the terms
with $\abs{\iota'}\in\{1,r-1\}$, one factor has order $r$ and the other order $2$;
the remaining terms involve only orders between $3$ and $r-1$. By
\cref{lem:flow-linear}, applied as before with the last sum as the source, we
have
\[
 \Theta_r(\tau)\le\norm{D^rf}_\infty+C_{\ell,r}b\int_\tau^N\lambda(\tau')
 \abs{B(\tau')}\Bigl\{\Theta_2\Theta_r+\sum_{q=3}^{r-1}\Theta_q\Theta_{r+2-q}\Bigr\}
 (\tau')\dd\tau'.
\]
By (b) and induction on $r$, the sum over $q$ is bounded. Gronwall's inequality
then bounds $\Theta_r$, with a constant of the stated type. This proves (c).
\end{proof}

\begin{proof}[Proof of \cref{lem:flow-step-limit}]
We begin with (a). The gradient bound is \cref{lem:flow-spatial}(a). By
\cref{lem:inp-steps}(a), each integration step is monotone and commutes with
the addition of constants, which implies that it does not increase the supremum
of the difference of two functions. Composing the steps, we find that
$\sup\abs{F-F'}\le\sup\abs{f-f'}$. For \eqref{eq:flow-L1}, we refine the partitions
so that both $m_j$ and $\widetilde m_j$ are constant on each piece; by
\cref{lem:inp-steps}(b) this does not change $F$ or $\widetilde F$. Because $B_j$
is symmetric, we have
\[
 m_j\langle\nabla F,B_j\nabla F\rangle-\widetilde m_j\langle\nabla\widetilde F,
 B_j\nabla\widetilde F\rangle
 =m_j\bigl\langle B_j(\nabla F+\nabla\widetilde F),\nabla(F-\widetilde F)\bigr\rangle
 +(m_j-\widetilde m_j)\langle\nabla\widetilde F,B_j\nabla\widetilde F\rangle.
\]
Then $\mathcal D=F-\widetilde F$ satisfies, on stage $j$,
\[
 -\partial_\tau \mathcal D=\frac{b\lambda_j}{2}B_j\mathbin{:}D^2\mathcal D
 +\frac{b\lambda_jm_j}{2}\bigl\langle B_j(\nabla F+\nabla\widetilde F),\nabla \mathcal D\bigr\rangle
 +\frac{b\lambda_j}{2}(m_j-\widetilde m_j)\langle\nabla\widetilde F,B_j\nabla\widetilde F\rangle,
\]
with $\mathcal D(N,\cdot)=0$. By \cref{lem:flow-cole-hopf}(a) and (b), $\mathcal D$ is of class
$C^{1,2}$ on each piece, $\nabla \mathcal D$ is bounded, and the drift is bounded and
Lipschitz in $x$, since the Hessians of $F$ and $\widetilde F$ are bounded on each
piece. The source is bounded by $(b/2)\lambda_j\abs{B_j}\abs{m_j-\widetilde m_j}L_f^2$,
by the gradient bound. Then \cref{lem:flow-linear} implies \eqref{eq:flow-L1}.

For (b), by \eqref{eq:flow-L1}, we see that the step flows $F^{(l)}$ of any
sequence with $\int_0^1\abs{m_j^{(l)}-m_j}\to0$ form a Cauchy sequence for
uniform convergence. Interlacing two such sequences gives another one, and we
conclude that the limit does not depend on the sequence. The estimate
\eqref{eq:flow-L1} and the bound on $\sup\abs{F-F'}$ pass to uniform limits. For the
gradient, let $e_{l,l'}=\sup_{\tau,x}\abs{F^{(l)}-F^{(l')}}$, let $v$ be a unit
vector, and let $a>0$. By Taylor's formula and the Hessian bound of
\cref{lem:flow-spatial}(a), which holds for every $l$, the difference quotient
$(F^{(l)}(\tau,x+av)-F^{(l)}(\tau,x))/a$ differs from $\partial_vF^{(l)}(\tau,x)$ by
at most $(K_f+L_f^2)a/2$. Comparing the difference quotients of $F^{(l)}$ and
$F^{(l')}$, we find that
\[
 \sup_{\tau,x}\abs{\partial_vF^{(l)}(\tau,x)-\partial_vF^{(l')}(\tau,x)}
 \le\frac{2e_{l,l'}}a+(K_f+L_f^2)a.
\]
Choosing $a=e_{l,l'}^{1/2}$, we see that the gradients of the step flows are
uniformly Cauchy. Their uniform limit is continuous, and by the fundamental
theorem of calculus it is the gradient of $F$, which therefore satisfies
$\abs{\nabla F}\le L_f$.

For (c), let $F^{(l)}$ be step flows of stages $1,\dots,N$ with masses
$m_j^{(l)}$ as in (b), and let $f_{\rm in}^{(l)}=F^{(l)}(N',\cdot)$, which is the
output of the step flow of stages $N'+1,\dots,N$. Let $\bar F^{(l)}$ and $\bar F$ be the
step flow of stages $1,\dots,N'$ with masses $m_j^{(l)}$ and terminal function
$f_{\rm in}$, and its limit. On $[0,N']$, the function $F^{(l)}$ is the step flow
of stages $1,\dots,N'$ with terminal function $f^{(l)}_{\rm in}$. By part (a), we have
\[
 \sup\abs{F^{(l)}-\bar F}\le\sup\abs{f^{(l)}_{\rm in}-f_{\rm in}}+\sup\abs{\bar F^{(l)}-\bar F}.
\]
Both terms tend to zero, by (b) for the stages $N'+1,\dots,N$ and for the stages
$1,\dots,N'$. This proves (c).

For (d), we first note that the hypothesis of \cref{lem:flow-spatial}(b) does
not involve the masses and holds for every step flow. Using
\cref{lem:flow-spatial}, we see that all spatial derivatives of the step flows are
bounded by constants of the type stated there, uniformly in $l$. Let $l,l'$ be
two indices, refine the partitions so that $m_j^{(l)}$ and $m_j^{(l')}$ are
constant on each piece, and let $\mathcal D=F^{(l)}-F^{(l')}$. Then $\mathcal D$
satisfies the equation in the proof of (a), with $F=F^{(l)}$,
$\widetilde F=F^{(l')}$, and the masses $m_j^{(l)}$, $m_j^{(l')}$ in place of
$m_j$, $\widetilde m_j$. For $\iota\ne0$, we apply $D^\iota$ to that equation and
use the Leibniz rule. We find that
\[
 -\partial_\tau D^\iota\mathcal D=\frac{b\lambda_j}{2}B_j\mathbin{:}D^2D^\iota\mathcal D
 +\frac{b\lambda_jm^{(l)}_j}{2}\bigl\langle B_j(\nabla F^{(l)}+\nabla F^{(l')}),
 \nabla D^\iota\mathcal D\bigr\rangle+\Upsilon_\iota,
\]
where $\Upsilon_\iota$ is the sum of the terms
$(b\lambda_jm^{(l)}_j/2)\binom\iota{\iota'}\langle B_jD^{\iota'}\nabla(F^{(l)}+F^{(l')}),
\nabla D^{\iota-\iota'}\mathcal D\rangle$ with $0\ne\iota'\le\iota$ and of
$(b\lambda_j/2)(m^{(l)}_j-m^{(l')}_j)D^\iota\langle\nabla F^{(l')},B_j\nabla F^{(l')}\rangle$.
The terms with $\abs{\iota'}=1$ are linear in the derivatives of $\mathcal D$ of
order $\abs\iota$, with coefficients bounded by $C\lambda_j\abs{B_j}$; the terms
with $\abs{\iota'}\ge2$ involve derivatives of $\mathcal D$ of order less than
$\abs\iota$; and the last term is bounded by
$C\lambda_j\abs{B_j}\abs{m^{(l)}_j-m^{(l')}_j}$, with $C$ depending on the
derivatives of $F^{(l')}$ of order at most $\abs\iota+1$. By
\cref{lem:flow-linear}, induction on $\abs\iota$, and Gronwall's inequality, we
have
\[
 \sup_{\tau,x}\abs{D^\iota\mathcal D(\tau,x)}
 \le C\sum_j\lambda_j\int_0^1\abs{m^{(l)}_j-m^{(l')}_j}\dd\tau.
\]
It follows that every spatial derivative of the step flows converges uniformly,
and the limit $F$ has continuous spatial derivatives of all orders, which are the
limits of those of the step flows. Letting $l'\to\infty$ in the last
display gives \eqref{eq:flow-L1-derivatives}, and the Hessian bound follows from
\cref{lem:flow-spatial}(b).

Now let each $m_j$ be continuous, and choose the step flows with the sampled
masses $m_j^{\mathfrak p_l}$ for partitions with $\abs{\mathfrak p_l}\to0$; then
$m_j^{\mathfrak p_l}\to m_j$ pointwise and in $L^1$. The step flows are
continuous in $\tau$ and continuously differentiable on each piece, and, for
$j-1\le\tau\le\tau_1\le j$, they satisfy
\begin{multline*}
 F^{(l)}(\tau,x)=F^{(l)}(\tau_1,x)\\
 +\int_\tau^{\tau_1}\frac{b\lambda_j}{2}
 \Bigl\{B_j\mathbin{:}D^2F^{(l)}+m_j^{\mathfrak p_l}(\tau'-j+1)
 \bigl\langle\nabla F^{(l)},B_j\nabla F^{(l)}\bigr\rangle\Bigr\}(\tau',x)\dd\tau'.
\end{multline*}
By bounded convergence, the same identity holds for $F$ and $m_j$. Its integrand
is continuous in $\tau'$, since $\nabla F$ and $D^2F$ are uniform limits of
continuous functions and $m_j$ is continuous. This shows that $F$ is
continuously differentiable in $\tau$ on $[j-1,j]$ and solves
\eqref{eq:flow-equation}. For uniqueness, let $F'$ be another solution with the
stated properties. The difference $\mathcal D=F-F'$ satisfies the equation in
the proof of (a), with $\widetilde F=F'$, $\widetilde m_j=m_j$, and
$\mathcal D(N,\cdot)=0$. Its drift is bounded and Lipschitz in $x$, since the first
and second derivatives of $F$ and $F'$ are bounded, and $\mathcal D$ and
$\nabla\mathcal D$ have linear growth.
Then \cref{lem:flow-linear}, with $\Upsilon=0$, implies that $\mathcal D=0$.
\end{proof}

\section{Parameter derivatives}\label{sec:flow-parameters}

We now prove \cref{prop:flow-uniform}. Throughout this section we assume
\ref{H1}--\ref{H5}, and on stage $j$ we abbreviate
\[
 \mathcal A_j=\frac{bh}{2}\kappa_jB_j,\qquad
 \mathcal V_j=\frac{bh}{2}\kappa_jm_jB_j,
\]
so that \eqref{eq:flow-equation} reads
$-\partial_\tau F=\mathcal A_j\mathbin{:}D^2F+\langle\nabla F,\mathcal V_j\nabla F\rangle$.
Both matrices are symmetric and depend on $(h,k)$, and $\mathcal V_j$ also depends
on $\tau$ through the mass. Their derivatives
$\partial^{\gamma'}\mathcal A_j=(b/2)\partial^{\gamma'}(h\kappa_jB_j)$ and
$\partial^{\gamma'}\mathcal V_j$ are bounded by the suprema listed in
\cref{prop:flow-uniform}(ii). Note that $\partial_h\mathcal A_j$ does not vanish
at $h=0$; these derivatives are bounded but not small. The proof has two
parts. In \cref{lem:flow-step-bounds}, we prove the bounds for the step flows
with sampled masses, with constants that do not depend on the partitions. We
then pass to the limit.

\begin{lemma}[Uniform bounds for step flows]\label{lem:flow-step-bounds}
Let $\mathfrak p$ be partitions, and let $F^{\mathfrak p}(\tau,h,k,x)$ be the step
flow with the sampled masses $m_j^{\mathfrak p}(\cdot,h,k)$ and terminal function
$f(h,k,\cdot)$, for $(h,k)\in\mathcal U_0$. Then $F^{\mathfrak p}$ is infinitely
differentiable in $(h,k,x)$, its derivatives $\partial^\gamma D^\iota F^{\mathfrak p}$
are continuous in $(\tau,h,k,x)$ and, on each closed piece, continuously
differentiable in $\tau$, with $\partial_\tau$ commuting with $\partial^\gamma
D^\iota$. For every $(\gamma,\iota)\ne(0,0)$, the function
$\partial^\gamma D^\iota F^{\mathfrak p}$ is bounded on $[0,N]\times\mathcal
U_0\times\R^\ell$ by a constant $C$ of the type stated in
\cref{prop:flow-uniform}(ii), which does not depend on $\mathfrak p$.
\end{lemma}

\begin{proof}
On the piece $(j,i)$, the step \eqref{eq:flow-step-solution} has mass
$m_j(\tau_{j,i-1},h,k)$ and covariance $t\,bh\kappa_j(h,k)B_j(h,k)$, with
$t=j-1+\tau_{j,i}-\tau\in[0,1]$. Both are smooth functions of $(h,k)$ on the
compact convex set $\mathcal U_0$, and the covariance is positive semidefinite
for $h\ge0$. Applying \cref{lem:flow-cole-hopf}(c) to the pieces one after
another, from the innermost, we obtain the regularity statements. The
commutation of $\partial_\tau$ with $\partial^\gamma D^\iota$ holds because
$F^{\mathfrak p}$ is infinitely differentiable in $(\tau,h,k,x)$ on each closed
piece. It remains to prove the bounds, which we do by induction. We drop the
superscript $\mathfrak p$ and write $F_\gamma=\partial^\gamma F$ and
\[
 \Theta_{\gamma,r}(\tau)=\max_{\abs\iota=r}\sup_{x}
 \abs{D^\iota F_\gamma(\tau,h,k,x)},
\]
for fixed $(h,k)\in\mathcal U_0$; all bounds below are uniform in $(h,k)$.

\paragraph{Step 1: spatial derivatives.} For $\gamma=0$ and $r\ge1$, the bound
on $\Theta_{0,r}$ is \cref{lem:flow-spatial}, which applies since
$2\ell b\norm{D^2f}_\infty\sum_j\lambda_j\abs{B_j}\le2\ell b\,h_0\Lambda_*
\norm{D^2f}_\infty\le1$ by \ref{H5}. We have $\Theta_{0,2}\le2\norm{D^2f}_\infty$.

\paragraph{Step 2: the differentiated equations.} Let $\gamma\ne0$. Applying
$\partial^\gamma$ to \eqref{eq:flow-equation} and separating the terms that
contain $F_\gamma$, we find that, on each piece of stage $j$,
\begin{equation}\label{eq:flow-Fgamma}
 -\partial_\tau F_\gamma=\mathcal A_j\mathbin{:}D^2F_\gamma
 +2\langle\mathcal V_j\nabla F,\nabla F_\gamma\rangle+\Upsilon_\gamma,
\end{equation}
with
\begin{equation}\label{eq:flow-Ugamma}
 \Upsilon_\gamma=\sum_{0\ne\gamma'\le\gamma}\binom{\gamma}{\gamma'}
 \partial^{\gamma'}\!\mathcal A_j\mathbin{:}D^2F_{\gamma-\gamma'}
 +\sum_{\substack{\gamma'+\gamma''+\gamma'''=\gamma\\ \gamma''\ne\gamma,\ \gamma'''\ne\gamma}}
 \frac{\gamma!}{\gamma'!\,\gamma''!\,\gamma'''!}
 \bigl\langle\nabla F_{\gamma''},\partial^{\gamma'}\!\mathcal V_j\nabla F_{\gamma'''}\bigr\rangle,
\end{equation}
where the second sum runs over the triples of multi-indices with
$\gamma'+\gamma''+\gamma'''=\gamma$ other than $(0,\gamma,0)$ and $(0,0,\gamma)$; since
$\gamma\ne0$, these are the triples with $\gamma''\ne\gamma$ and $\gamma'''\ne\gamma$.
Every $F_{\gamma''}$, $F_{\gamma'''}$, and $F_{\gamma-\gamma'}$ appearing in
$\Upsilon_\gamma$ has parameter order strictly less than $\abs\gamma$, and
$F_0=F$ enters only through $\nabla F$ and $D^2F$. Applying $D^\iota$ with
$\abs\iota=r$ to \eqref{eq:flow-Fgamma}, we obtain
\begin{multline}\label{eq:flow-Fgamma-iota}
 -\partial_\tau D^\iota F_\gamma=\mathcal A_j\mathbin{:}D^2D^\iota F_\gamma
 +2\langle\mathcal V_j\nabla F,\nabla D^\iota F_\gamma\rangle\\
 +2\sum_{0\ne\iota'\le\iota}\binom{\iota}{\iota'}
 \bigl\langle\mathcal V_j\nabla D^{\iota'}F,\nabla D^{\iota-\iota'}F_\gamma\bigr\rangle
 +D^\iota\Upsilon_\gamma,
\end{multline}
with terminal value $D^\iota F_\gamma(N,\cdot)=\partial^\gamma D^\iota f$.

\paragraph{Step 3: derivative bounds by weight.} We prove the bound on
$\Theta_{\gamma,r}$ by induction on $\abs\gamma$, and, for fixed $\gamma$, by
induction on $r=0,1,2,\dots$. \Cref{tab:flow-weights} lists the terms of
\eqref{eq:flow-Fgamma-iota} and the derivatives of $F$ that they contain. Let
$\varpi=r+2\abs\gamma$ be the weight of $D^\iota F_\gamma$. We read off from the
table that every derivative $\partial^{\gamma^\circ}D^{\iota^\circ}F$ other than
the unknown $D^\iota F_\gamma$ has $\gamma^\circ\le\gamma$ and weight
$\abs{\iota^\circ}+2\abs{\gamma^\circ}\le\varpi$. Moreover, either
$\abs{\gamma^\circ}<\abs\gamma$, or $\gamma^\circ=\gamma$ and $\abs{\iota^\circ}\le r$.
The terms with $\gamma^\circ=\gamma$ and $\abs{\iota^\circ}=r$ are those in the
second row. They are linear in the derivatives of order $r$ of $F_\gamma$, with
coefficients bounded by $2\ell r\abs{\mathcal V_j}\Theta_{0,2}$. In the fourth
row, a parameter derivative is paid for by two spatial derivatives. Indeed, when
$\partial^{\gamma'}$ falls on the coefficient $\mathcal A_j$ of $D^2$, the
derivative $F_{\gamma-\gamma'}$ appears with $r+2$ spatial derivatives, and its
weight $r+2+2(\abs\gamma-\abs{\gamma'})$ is at most $\varpi$ because
$\abs{\gamma'}\ge1$.

\begin{table}[ht]
\centering
\small
\caption{The terms of \eqref{eq:flow-Fgamma-iota}, for $\abs\iota=r$ and
$\varpi=r+2\abs\gamma$. The second column lists the derivatives of $F$ that each term
contains, and the third says how they are controlled.}\label{tab:flow-weights}
\begin{tabular}{@{}>{\raggedright\arraybackslash}p{0.33\textwidth}>{\raggedright\arraybackslash}p{0.36\textwidth}>{\raggedright\arraybackslash}p{0.23\textwidth}@{}}
\toprule
Term & Derivatives of $F$ & Control\\
\midrule
$\mathcal A_j\mathbin{:}D^2D^\iota F_\gamma$,
$2\langle\mathcal V_j\nabla F,\nabla D^\iota F_\gamma\rangle$
 & the unknown $D^\iota F_\gamma$ & operator of \cref{lem:flow-linear}\\
$\langle\mathcal V_j\nabla D^{\iota'}F,\nabla D^{\iota-\iota'}F_\gamma\rangle$,
$\abs{\iota'}=1$
 & $D^2F$; derivatives of order $r$ of $F_\gamma$ (weight $\varpi$)
 & Gronwall's inequality\\
the same, $\abs{\iota'}\ge2$
 & $D^{\abs{\iota'}+1}F$ (weight $\le r+1$); derivatives of order
 $r-\abs{\iota'}+1<r$ of $F_\gamma$ & induction on $r$\\
$\partial^{\gamma'}\mathcal A_j\mathbin{:}D^2D^\iota F_{\gamma-\gamma'}$,
$\gamma'\ne0$
 & order $r+2$ of $F_{\gamma-\gamma'}$, weight $\le\varpi$ & induction on $\abs\gamma$\\
$\langle D^{\iota_1}\nabla F_{\gamma''},\partial^{\gamma'}\mathcal V_j
D^{\iota_2}\nabla F_{\gamma'''}\rangle$, $\iota_1+\iota_2=\iota$
 & orders $\abs{\iota_1}+1$, $\abs{\iota_2}+1\le r+1$ of $F_{\gamma''}$,
 $F_{\gamma'''}$ with $\abs{\gamma''},\abs{\gamma'''}<\abs\gamma$; weights $<\varpi$
 & induction on $\abs\gamma$\\
\bottomrule
\end{tabular}
\end{table}

In the last row, $\gamma''=0$ or $\gamma'''=0$ is allowed; the corresponding
factor is a spatial derivative of $F$ of order at most $r+1\le\varpi-1$, bounded in
Step~1. By the induction hypothesis and Step~1, we conclude that all terms of
\eqref{eq:flow-Fgamma-iota} other than those in the first two rows are bounded by
a constant $c_0$. The coefficients in the second row are bounded by
$c_1=2\ell r\max_j\sup\abs{\mathcal V_j}\cdot2\norm{D^2f}_\infty$.

\paragraph{Step 4: the estimate.} We apply \cref{lem:flow-linear} to each
$\psi=D^\iota F_\gamma$ with $\abs\iota=r$, with $\mathcal A=\mathcal A_j$ on stage
$j$, the drift $\phi=2\mathcal V_j\nabla F$, and the source formed by all other
terms of \eqref{eq:flow-Fgamma-iota}. The drift is bounded by Step~1, and it is
Lipschitz in $x$ since $D^2F$ is bounded. By the regularity statement, $\psi$,
$\nabla\psi$, and the source are bounded and continuous on each piece. Taking the
maximum over $\abs\iota=r$, we find
\[
 \Theta_{\gamma,r}(\tau)\le\max_{\abs\iota=r}\sup_x\abs{\partial^\gamma D^\iota f}
 +\int_\tau^N\bigl\{c_1\Theta_{\gamma,r}(\tau')+c_0\bigr\}\dd\tau',
\]
and by Gronwall's inequality we have
$\Theta_{\gamma,r}\le(\max_{\abs\iota=r}\sup\abs{\partial^\gamma D^\iota f}+Nc_0)e^{Nc_1}$.
For $r=0$ the second row is absent and $c_1=0$. This completes the induction.

The constant obtained for $(\gamma,r)$ involves the bounds for the pairs
$(\gamma^\circ,r^\circ)$ listed in \cref{tab:flow-weights}, the bounds on
$\partial^{\gamma'}\mathcal A_j$ and $\partial^{\gamma'}\mathcal V_j$ for
$\gamma'\le\gamma$, and the terminal values. Following the induction back, it
depends only on $\ell$, $N$, $b$, $h_1$, $\gamma$, $r$, and the suprema of
$\abs{\partial^{\gamma'}(\kappa_jB_j)}$, $\abs{\partial^{\gamma'}(\kappa_jm_jB_j)}$, and
$\abs{\partial^{\gamma'}D^{\iota'}f}$ over $\gamma'\le\gamma$ and
$0<\abs{\iota'}+2\abs{\gamma'}\le\varpi$, as claimed. None of them depends on
$\mathfrak p$, since the sampled masses have the same values and
$(h,k)$-derivatives as $m_j$ at the sampling points.
\end{proof}

\begin{proof}[Proof of \cref{prop:flow-uniform}]
Let $\mathfrak p_l$ be partitions with $\abs{\mathfrak p_l}\to0$, and let $F^{\mathfrak p_l}$
be the step flows with the sampled masses. The proof has three steps: convergence
of the values, convergence of the derivatives, and the equations.

\paragraph{Step 1: values.} By \ref{H3}, $m_j$ is uniformly continuous on
$[0,1]\times\mathcal U$, and we have
$\varepsilon_l=\sup_{\mathcal U}\sum_j\int_0^1\abs{m_j^{\mathfrak p_l}-m_j}\dd\tau\to0$.
Then \cref{lem:flow-step-limit}(a) and (b), applied at each $(h,k)$ with
$\sum_j\lambda_j\abs{B_j}\le h_0\Lambda_*$, give
\[
 \sup_{[0,N]\times\mathcal U_0\times\R^\ell}\abs{F^{\mathfrak p_l}-F}
 \le\frac b2\sup\abs{\nabla_xf}^2\,h_0\Lambda_*\,\varepsilon_l\longrightarrow0.
\]

\paragraph{Step 2: derivatives.} Fix $(\gamma,\iota)\ne(0,0)$. By
\cref{lem:flow-step-bounds}, the functions $\partial^\gamma D^\iota F^{\mathfrak p_l}$ are
bounded uniformly in $l$, and so are their first derivatives in $x$ and in
$(h,k)$, whose weights are larger by one and by two, respectively. Their derivative in
$\tau$ is given on each piece by \eqref{eq:flow-Fgamma-iota}, and it is also bounded
uniformly in $l$, since it involves derivatives of weight at most
$\abs\iota+2\abs\gamma+2$. Since $\partial^\gamma D^\iota F^{\mathfrak p_l}$ is also continuous in
$\tau$, the family $\{\partial^\gamma D^\iota F^{\mathfrak p_l}\}_l$ is uniformly bounded and uniformly
Lipschitz on $[0,N]\times\mathcal U_0\times\R^\ell$. By the Arzel\`a--Ascoli
theorem and a diagonal argument over the compact sets
$[0,N]\times\mathcal U_0\times\{\abs x\le R\}$, every subsequence has a further
subsequence along which, for all $(\gamma,\iota)$ at once,
$\partial^\gamma D^\iota F^{\mathfrak p_l}$ converges uniformly on compact sets to a continuous
limit $L_{\gamma,\iota}$. We claim that $L_{\gamma,\iota}=\partial^\gamma
D^\iota F$, so that the limits do not depend on the subsequence and the
whole sequence converges. Indeed, $L_{0,0}=F$ by Step~1. If $\varphi_l\to\varphi$ and
$\partial_a\varphi_l\to\psi$ uniformly on compact sets, where $\partial_a$ is a
derivative in $h$, in $k$, or in a coordinate of $x$, then letting
$l\to\infty$ in $\varphi_l(y+\eta e_a)-\varphi_l(y)=\int_0^\eta\partial_a\varphi_l(y+\eta'e_a)\dd\eta'$
shows that $\partial_a\varphi=\psi$; at boundary points of $\mathcal U_0$, such as $h=0$,
we take $\eta$ of the sign that keeps $y+\eta e_a$ in the domain. Induction on $\abs\gamma+\abs\iota$
proves the claim. This proves (iv), the existence and continuity of all
derivatives in (ii), and the bounds in (ii), which pass to the limit. For the
linear growth, note that $\abs{F(\tau,h,k,x)-F(\tau,h,k,0)}\le\sup\abs{\nabla_xf}
\abs x$, that $\abs{F(\tau,h,k,0)-f(h,k,0)}\le N\sup\abs{\partial_\tau F}$ is
bounded by \eqref{eq:flow-equation} and the bounds on $\nabla F$ and $D^2F$, and
that $\sup_{\mathcal U}\abs{f(h,k,0)}$ is finite, since $f(\cdot,\cdot,0)$ is
continuous on the compact set $\mathcal U$. The constant obtained depends on this
supremum in addition to the derivative bounds.

\paragraph{Step 3: the equations.} Part (i) is \cref{lem:flow-step-limit}(d),
whose hypothesis holds at each $(h,k)\in\mathcal U_0$ by \ref{H5}. For (iii), on
stage $j$ and for $j-1\le\tau\le\tau_1\le j$, the step flows satisfy
\[
 F^{\mathfrak p_l}_\gamma(\tau)=F^{\mathfrak p_l}_\gamma(\tau_1)+\int_\tau^{\tau_1}
 \bigl\{\mathcal A_j\mathbin{:}D^2F^{\mathfrak p_l}_\gamma
 +2\langle\mathcal V_j^{\mathfrak p_l}\nabla F^{\mathfrak p_l},\nabla F^{\mathfrak p_l}_\gamma\rangle
 +\Upsilon^{\mathfrak p_l}_\gamma\bigr\}(\tau')\dd\tau',
\]
where $F^{\mathfrak p_l}_\gamma=\partial^\gamma F^{\mathfrak p_l}$ and where $\mathcal V_j^{\mathfrak p_l}$ and
$\Upsilon^{\mathfrak p_l}_\gamma$ are formed with the sampled masses. By \ref{H3}, the
$(h,k)$-derivatives of $m_j^{\mathfrak p_l}$ converge pointwise to those of $m_j$,
and by (iv) the derivatives of $F^{\mathfrak p_l}$ converge pointwise. All terms are bounded,
and by bounded convergence the same identity holds for $F_\gamma=\partial^\gamma F$,
with the masses $m_j$. Since its integrand is continuous in $\tau'$, the function
$F_\gamma$ is continuously differentiable in $\tau$ on $[j-1,j]$ and satisfies
\eqref{eq:flow-Fgamma}, which is the equation obtained by applying $\partial^\gamma$
to \eqref{eq:flow-equation}. This proves (iii).
\end{proof}

\section{Response to a change of the terminal function}\label{sec:flow-response}

We now compare a hierarchy with terminal function $f$ and the same hierarchy with
terminal function $f+V$, for a bounded function $V$. The first-order change is
the expectation of $V$ under the path law of the unperturbed hierarchy, and the
error is quadratic in $\norm V_\infty$, with a constant that does not depend on
the number of levels. For a hierarchy with $M$ levels, let $\Sigma_1,\dots,\Sigma_M
\succeq0$ be the covariances and $m_1,\dots,m_M\in[0,1]$ the masses, listed from
the root; the masses need not be ordered. For a function $f$ of linear growth, the
partial values are $\Psi^f_M=f$ and $\Psi^f_{j-1}=T_{m_j,\Sigma_j}\Psi^f_j$ for
$j=M,\dots,1$, as in \eqref{eq:inp-hierarchy-recursion}, and we write
$\mathcal Hf=\Psi_0^f$. For $z_0\in\R^\ell$, we let $\E_f$ denote expectation for
the Markov chain $z_0,z_1,\dots,z_M$ in which $z_j$ has law
$\mathcal K_{m_j,\Sigma_j,\Psi^f_j}(z_{j-1},\cdot)$ given $z_0,\dots,z_{j-1}$.
This is the path law of \cref{def:inp-hierarchy}(c), started at $z_0$ and without
the terminal draw of spins.

\begin{lemma}[Response of a hierarchy to its terminal function]\label{b:lem:os-prefix}
Let $f$ have linear growth, and let $V\colon\R^\ell\to\R$ be bounded and
measurable. Then, for every $z_0\in\R^\ell$,
\[
 \bigl\lvert\mathcal H(f+V)(z_0)-\mathcal Hf(z_0)-\E_f V(z_M)\bigr\rvert
 \le\frac12\norm V_\infty^2.
\]
The bound does not depend on the number $M$ of levels or on the covariances.
\end{lemma}

The argument is the martingale argument that the proof of
\cite[Lemma~3.1]{oai-barriers} uses to show that $\chi\le1-u^2$, and the
differentiation of a hierarchy behind \cref{inp:ob-hierarchy-derivative}, namely
\cite[Lemma~2.2]{oai-barriers}; we apply them to the terminal function instead
of the field.

\begin{proof}
For $0\le t\le1$, let $\Psi^t_j=\Psi^{f+tV}_j$ be the partial values with
terminal function $f+tV$. Using the monotonicity and the translation identity of
\cref{lem:inp-steps}(a), we have $\abs{\Psi_j^t-\Psi_j^0}\le t\norm V_\infty$, and
all integrals below are finite. Let $\E^t_j$ and $\Var_j^t$ denote the expectation and
the variance under the kernel $\mathcal K_{m_j,\Sigma_j,\Psi^t_j}(z,\cdot)$ for a
fixed starting point $z$. We claim that $t\mapsto\Psi^t_{j-1}(z)$ is twice
continuously differentiable, with
\begin{equation}\label{eq:flow-response-recursion}
 \partial_t\Psi^t_{j-1}(z)=\E^t_j\bigl[\partial_t\Psi^t_j\bigr],\qquad
 \partial_t^2\Psi^t_{j-1}(z)=\E^t_j\bigl[\partial_t^2\Psi^t_j\bigr]
 +m_j\Var^t_j\bigl(\partial_t\Psi^t_j\bigr),
\end{equation}
where the expectation and the variance act on the function $y\mapsto\partial_t\Psi^t_j(y)$
or $\partial^2_t\Psi^t_j(y)$. Indeed, let $\mathcal Z_t$ be random variables, twice
differentiable in $t$ with $\abs{\partial_t\mathcal Z_t}+\abs{\partial^2_t\mathcal Z_t}$ bounded and
$\E e^{\abs{\mathcal Z_0}}<\infty$, and let $\mathfrak a(t)=m^{-1}\log\E e^{m\mathcal Z_t}$ for some
$m\in(0,1]$. Differentiation under the expectation gives
$\mathfrak a'=\E[\partial_t\mathcal Z_t\,e^{m\mathcal Z_t}]/\E e^{m\mathcal Z_t}$ and
\[
 \mathfrak a''=\frac{\E[\partial_t^2\mathcal Z_t\,e^{m\mathcal Z_t}]}{\E e^{m\mathcal Z_t}}
 +m\Bigl(\frac{\E[(\partial_t\mathcal Z_t)^2\,e^{m\mathcal Z_t}]}{\E e^{m\mathcal Z_t}}
 -\Bigl(\frac{\E[\partial_t\mathcal Z_t\,e^{m\mathcal Z_t}]}{\E e^{m\mathcal Z_t}}\Bigr)^2\Bigr).
\]
For $m=0$ and $\mathfrak a(t)=\E\mathcal Z_t$, the same formulas hold without the variance term. We
apply this with $\mathcal Z_t=\Psi^t_j(z+\Sigma_j^{1/2}\mathbf g)$ and $m=m_j$; the weights
$e^{m_j\mathcal Z_t}/\E e^{m_j\mathcal Z_t}$ are the densities of the kernel
$\mathcal K_{m_j,\Sigma_j,\Psi^t_j}(z,\cdot)$. Starting from $\partial_t\Psi^t_M=V$
and $\partial^2_t\Psi^t_M=0$, induction on $M-j$ shows that the derivatives are
bounded and proves \eqref{eq:flow-response-recursion}.

Now let $z_0,\dots,z_M$ be the chain of the lemma for the terminal function $f+tV$,
and let $\mathcal M_j=\partial_t\Psi^t_j(z_j)$. By the first identity in
\eqref{eq:flow-response-recursion} and the Markov property, we have
$\E^t[\mathcal M_j\mid z_0,\dots,z_{j-1}]=\mathcal M_{j-1}$. This implies that
$(\mathcal M_j)$ is a martingale, with $\mathcal M_M=V(z_M)$ and
$\mathcal M_0=\partial_t\mathcal H(f+tV)(z_0)$. Iterating the second identity in
\eqref{eq:flow-response-recursion} from $\partial^2_t\Psi^t_M=0$, we obtain
\[
 \partial_t^2\mathcal H(f+tV)(z_0)=\sum_{j=1}^Mm_j\,\E^t\bigl(\mathcal M_j-\mathcal M_{j-1}\bigr)^2.
\]
Since $0\le m_j\le1$ and the increments of a martingale are orthogonal,
\[
 0\le\partial_t^2\mathcal H(f+tV)(z_0)\le\sum_{j=1}^M\E^t\bigl(\mathcal M_j-\mathcal M_{j-1}\bigr)^2
 =\E^tV(z_M)^2-\bigl(\E^tV(z_M)\bigr)^2\le\norm V_\infty^2.
\]
At $t=0$, we have $\partial_t\mathcal H(f+tV)(z_0)=\mathcal M_0=\E_fV(z_M)$. By
Taylor's formula with integral remainder,
\[
 \mathcal H(f+V)(z_0)-\mathcal Hf(z_0)-\E_fV(z_M)
 =\int_0^1(1-t)\,\partial_t^2\mathcal H(f+tV)(z_0)\dd t,
\]
which lies in $[0,\norm V_\infty^2/2]$. This completes the proof.
\end{proof}

We next apply \cref{b:lem:os-prefix} to the scalar hierarchy that integrates a
block of $\mathfrak n$ identical coordinates over a time interval $[0,t_\star]$, and pass to
a general distribution function. For a distribution function $\alpha$, an
integer $\mathfrak n\ge1$, a time $t_\star\in(0,1]$, and $f\in\mathcal C^2(\R)$, let
$\mathcal H^\alpha_{t_\star,\mathfrak n}f$ denote the value at $z=0$ of the flow of
\cref{lem:flow-step-limit}(b) with $\ell=1$, $N=1$, $B_1=1$, $\lambda_1=t_\star$,
mass $m_1(\tau)=\alpha(t_\star\tau)/\mathfrak n$, and terminal function $f$. If
$\alpha_n$ are step distribution functions with $\int_0^1\abs{\alpha_n-\alpha}\dd t\to0$,
then $\mathcal H^\alpha_{t_\star,\mathfrak n}f$ is the limit of the values at $z=0$ of the
step flows with masses $\alpha_n(t_\star\tau)/\mathfrak n$. In the scalar time
$t=t_\star\tau$, these step flows solve
$-\partial_tF=(b/2)\{\partial_{zz}F+(\alpha_n(t)/\mathfrak n)(\partial_zF)^2\}$ on each
interval where $\alpha_n$ is constant, with $F(t_\star,\cdot)=f$. For a general
$\alpha$, we do not assert that the limiting flow is differentiable in time at
the jumps of $\alpha$. Recall the function $\Phi=\Phi_\alpha$ and the diffusion
$Y$ of \eqref{b:eq:scalar-diffusion}.

\begin{corollary}[Response of the common part of a hierarchy]\label{cor:flow-prefix}
Let $\alpha$ be a distribution function, let $\mathfrak n\ge1$ be an integer, and let
$t_\star\in(0,1]$. Then $\mathcal H^\alpha_{t_\star,\mathfrak n}(\mathfrak n\Phi_\alpha(t_\star,\cdot))
=\mathfrak n\Phi_\alpha(0,0)$. Moreover, let $V\in C^2(\R)$ with $V$, $V'$, and $V''$ bounded,
and suppose that $\norm V_\infty\le\eps$. Then
\[
 \bigl\lvert\mathcal H^\alpha_{t_\star,\mathfrak n}\bigl(\mathfrak n\Phi_\alpha(t_\star,\cdot)+V\bigr)
 -\mathfrak n\Phi_\alpha(0,0)-\E V(Y_{t_\star})\bigr\rvert\le\frac{\eps^2}{2}.
\]
\end{corollary}

In \cref{sec:cubic}, we apply the corollary with $\mathfrak n=3$, with $t_\star$ the left
endpoint of the window, and with $V$ the difference between the restriction of
the local value to the diagonal and $3\Phi_\alpha(t_\star,\cdot)$, for which
$\eps=Ch^2k^2$. The initial segment of the hierarchy then contributes $\E V(Y_{t_\star})+O(h^4k^4)$. Note that
the law of $Y_{t_\star}$ does not depend on the perturbation.

\begin{proof}
Let $\alpha_n$ be step distribution functions with $t_\star$ among their breakpoints
and $\int_0^1\abs{\alpha_n-\alpha}\dd t\to0$; for instance, let $\alpha_n$ be
equal to $\alpha(t_{n,i-1})$ on $[t_{n,i-1},t_{n,i})$, where the $t_{n,i}$ are the
points $i2^{-n}$ together with $t_\star$. Let $\Phi_n=\Phi_{\alpha_n}$, let $Y^n$ be
the diffusion of $\alpha_n$, driven by the same Brownian motion as $Y$, and let
$\mathcal H_n$ denote the step flow with masses $\alpha_n(t_\star\tau)/\mathfrak n$, that is,
the finite hierarchy whose levels are the intervals between consecutive
breakpoints of $\alpha_n$ in $[0,t_\star]$.

\paragraph{Step 1: the step hierarchies.} By \eqref{eq:inp-scalar-recursion} and
\cref{lem:inp-steps}(d), the partial values of $\mathcal H_n$ with terminal
function $\mathfrak n\Phi_n(t_\star,\cdot)$ are $\mathfrak n\Phi_n(t_{n,i},\cdot)$, and in particular
$\mathcal H_n(\mathfrak n\Phi_n(t_\star,\cdot))(0)=\mathfrak n\Phi_n(0,0)$. The level from
$t_{n,i-1}$ to $t_{n,i}$ has mass $\alpha_n(t_{n,i-1})/\mathfrak n$ and partial value
$\mathfrak n\Phi_n(t_{n,i},\cdot)$. Its kernel has the density
$\exp\{\alpha_n(t_{n,i-1})(\Phi_n(t_{n,i},y)-\Phi_n(t_{n,i-1},z))\}$ with respect to
$N(z,b(t_{n,i}-t_{n,i-1}))$. This is the kernel of the scalar recursion of
$\alpha_n$ on the same interval.
By \cref{inp:ob-scalar}(iii), the path law of this hierarchy, started at $0$, is
the law of $(Y^n_{t_{n,i}})_{t_{n,i}\le t_\star}$. Then \cref{b:lem:os-prefix}, with
$f=\mathfrak n\Phi_n(t_\star,\cdot)$ and $z_0=0$, gives
\begin{equation}\label{eq:flow-prefix-step}
 \bigl\lvert\mathcal H_n\bigl(\mathfrak n\Phi_n(t_\star,\cdot)+V\bigr)(0)-\mathfrak n\Phi_n(0,0)
 -\E V(Y^n_{t_\star})\bigr\rvert\le\frac{\eps^2}{2}.
\end{equation}

\paragraph{Step 2: the limit.} By \cref{inp:ob-scalar}(i), $\abs{\partial_z\Phi_n}\le1$
and $0<\partial_{zz}\Phi_n\le1$, and the same holds for $\Phi_\alpha$; hence the
terminal functions below belong to $\mathcal C^2(\R)$, with gradients bounded by
$\mathfrak n+\norm{V'}_\infty$. Using \cref{lem:flow-step-limit}(a), we find that
\[
 \abs{\mathcal H_n(\mathfrak n\Phi_n(t_\star,\cdot)+V)-\mathcal H_n(\mathfrak n\Phi_\alpha(t_\star,\cdot)+V)}
 \le\mathfrak n\sup_z\abs{\Phi_n(t_\star,z)-\Phi_\alpha(t_\star,z)},
\]
which tends to zero by \cref{inp:ob-scalar}(ii). By \cref{lem:flow-step-limit}(b),
$\mathcal H_n(\mathfrak n\Phi_\alpha(t_\star,\cdot)+V)(0)\to\mathcal H^\alpha_{t_\star,\mathfrak n}
(\mathfrak n\Phi_\alpha(t_\star,\cdot)+V)$, since
$\int_0^1\abs{\alpha_n(t_\star\tau)-\alpha(t_\star\tau)}\dd\tau\to0$. Also,
\cref{inp:ob-scalar}(ii) implies that $\Phi_n(0,0)\to\Phi_\alpha(0,0)$. Finally, we write
$u_n=\partial_z\Phi_n$ and $u=\partial_z\Phi_\alpha$. Since $\abs{u_n}\le1$,
$\alpha\le1$, and $\abs{\partial_zu}\le1$, we have
\[
 \abs{\alpha_n(t)u_n(t,y')-\alpha(t)u(t,y)}
 \le\abs{\alpha_n(t)-\alpha(t)}+\sup_{t,z}\abs{u_n-u}+\abs{y'-y},
\]
and the equations for $Y^n$ and $Y$ give, for $t\le t_\star$,
\[
 \abs{Y^n_t-Y_t}\le b\int_0^1\abs{\alpha_n-\alpha}\dd t'+b\sup\abs{u_n-u}
 +b\int_0^t\abs{Y^n_{t'}-Y_{t'}}\dd t'.
\]
By Gronwall's inequality and \cref{inp:ob-scalar}(ii), $Y^n_{t_\star}\to Y_{t_\star}$
almost surely, and $\E V(Y^n_{t_\star})\to\E V(Y_{t_\star})$ by bounded convergence.
Taking $n\to\infty$ in \eqref{eq:flow-prefix-step} proves the second claim. The
first claim follows in the same way from the identity
$\mathcal H_n(\mathfrak n\Phi_n(t_\star,\cdot))(0)=\mathfrak n\Phi_n(0,0)$.
\end{proof}

\begin{remark}[A common block acts on the diagonal]\label{rem:flow-diagonal}
Let $e\in\R^\ell$, and suppose that stage $j$ of a stage datum has $B_j=ee^{\mathsf T}$.
Then the Gaussian increments of that stage are of the form $\zeta e$ with a
centered Gaussian variable $\zeta$, and a point of the line $\{ze:z\in\R\}$ stays
on it. It follows that, for $F\in\mathcal C^2(\R^\ell)$, the flow of this stage with
terminal function $F$, evaluated at $ze$, is the scalar flow with $B_j=1$, the
same $\lambda_j$ and $m_j$, and terminal function $z\mapsto F(ze)$, evaluated at
$z$. This holds for step flows by \cref{lem:inp-steps}(e) and passes to the limit
by \cref{lem:flow-step-limit}(b). In \cref{sec:cubic}, the part of the hierarchy
before the window has $B=ee^{\mathsf T}$ with $e=(1,1,1)$ and masses $\alpha/3$,
and its value is $\mathcal H^\alpha_{t_\star,3}$ applied to the restriction of the
local value to the diagonal.
\end{remark}
